\section{Introduction}

Language models can teach differently even when responding to the same exercise.
One reveals the answer; another asks the student to complete a step; a third
acknowledges frustration. Do these differences reflect recurring teaching
tendencies that predict behavior in a new educational situation? This is the
behavioral question about \emph{educational personality} examined here.

Distinguishing models' text is insufficient. Length, formatting, task
instructions, and solution access can all produce recognizable differences.
Conversely, a recurring tendency need not imply identical behavior in every
situation: a tutor may consistently elicit more reasoning after an erroneous
attempt. We therefore distinguish a model's default action profile from its
conditional response to student state, and test whether both add predictive
information beyond a baseline using educational context alone. Observed
differences, repeated behavior, and successful
transfer provide different kinds of evidence.

Educational actions make this question testable without assuming a human
personality taxonomy. Revealing a solution and inviting reasoning are observable,
can co-occur, and are not universally good or bad teaching moves. Acknowledging
expressed frustration does not establish empathy. Our target is the \emph{tutor
model's} behavior, distinct from adapting instruction to a \emph{student's}
assigned personality \citep{rooein2026pats}. Neither behavior predictability nor
its failure measures learning benefit.

Prior studies already investigate personality-oriented scenarios and situated
behavioral profiles \citep{lee2025trait,liu2026bma}, characteristic tutoring-action
mixtures, and learned tutor-persona variation
\citep{kucheria2025cima,lee2026tutorpersonas}. We examine a specific boundary in
these interpretations: does information that transfers to new problems also
transfer when students express themselves differently?

Three linked questions organize the study. \textbf{RQ1:} Under matched input,
do models show distinguishable defaults that recur across requests?
\textbf{RQ2:} Do defaults or model-specific student-state responses predict
new problems and new expressions? \textbf{RQ3:} How do neutral role rewordings
and explicit teaching instructions change those profiles?

An audit of a million-record educational-output archive motivates the controlled
study: different roles, repeated sources, and instructional settings can produce
apparent model differences (Appendix~\ref{app:archive-reanalysis}). We develop
observable measures on an excluded pilot, then fit profiles on training sources
and lock predictions before generating responses to new problems. Model-by-subject
baselines distinguish subject habits from student-state increments. After
observing this confirmation, a separately frozen expression extension reuses
its 32 problems and original forecasts, with contemporary original wording as
a control. The extension tests an expression shift without adding new sources.

Default revelation and elicitation gains survive both transfer tests. The
acknowledgement state profile improves new-problem prediction under original
wording, but loses its increment under explicit synonyms and implicit cues;
the contemporary original-wording control remains positive. Explicit teaching
instructions also make demanded actions converge while permissible accompanying
actions remain different. This separation of default recurrence, conditional
transfer, and instructional constraint is the central empirical contribution.
The finite constructed materials limit the supported interpretation to
prompt-dependent action profiles of the specified deployments.
